\section{The Pure Jump Libor Model}
Finally, in this section, we would like to provide a second example for the driving process $\Sigma$, so as to let the reader appreciate the degree of generality of this framework.
As in the general setup, we specify the process under the terminal probability measure $\mathbb{P}_{T_N}$. The example we choose is a matrix compound Poisson process, which was already presented in section \ref{pure_jump}:
\begin{align}
d\Sigma_t=M\Sigma_t+\Sigma_tM^{\top}+dL_t^{\mathbb{P}_{T_N}}.
\end{align}
All assumptions presented in section \ref{pure_jump} are in order. More specifically, we assume that $L_t^{\mathbb{P}_{T_N}}$ is a compound Poisson process with constant intensity $\lambda$ and jump distribution taking values in $S_d^{++}$. As a specific example of jump distribution we choose the standard Wishart distribution. By recalling the results in section \ref{pure_jump} we have that the solution for the function $\psi_\tau(u)$ is
\begin{align}
\psi_\tau(u)=e^{M^{\top}\tau}ue^{M\tau},
\end{align}
whereas for $\phi_\tau(u)$ we have
\begin{align}
\phi_\tau(u)=-\lambda\int_{0}^{\tau}{\det\left(I_d+2e^{M^{\top}s}ue^{Ms}\mathcal{Q}\right)^{-\frac{n}{2}}ds}+\lambda \tau.
\end{align}
In concrete pricing applications, the computation of the solution for $\phi_\tau(u)$ implies a numerical integration with respect to the time dimension. This numerical integration has an impact on the performance of the model which turns out to be slower than the Wishart Libor model. For illustrative purposes, we report an example for an implied volatility surface for caplets generated by the compound Poisson Libor model with central Wishart distributed jumps.
The mean reversion matrix $M$ and the jump intensity $\lambda$ are given by:
\begin{align}
M&=\left(\begin{array}{cc} -0.0550&0\\ 0& -0.1760\end{array}\right),\nonumber\\
\lambda&=0.1\nonumber.
\end{align}
As far as the jump size distribution is concerned, the parameters are the following:
\begin{align}
\mathcal{Q}&=\left(\begin{array}{cc} 0.27&0\\ 0& 0.05\end{array}\right),\nonumber\\
n&=3.1, \nonumber
\end{align}
and the initial state of the process is
\begin{align*}
\Sigma_0&=\left(\begin{array}{cc} 1.875&0.6\\ 0.6& 1.275\end{array}\right).
\end{align*}

\begin{center}
[Insert Figure \ref{fig:10}  here]
\end{center}
